%
\subsection*{Simetría y producto escalar}
\label{CH08-B01}
En un dominio simétrico, $f$ es par si $f(-x)=f(x)$ e impar si
$f(-x)=-f(x)$. Para funciones reales de cuadrado integrable, el producto
escalar y su norma son
\[
 \langle f,g\rangle=\int_{-L}^{L}f(x)g(x)\,dx,
 \qquad \|f\|^2=\langle f,f\rangle.
\]
Dos funciones son ortogonales cuando su producto escalar es cero.
Para $L>0$, la constante $1$, los cosenos y los senos de frecuencias
$n\pi/L$, $n\ge1$, son mutuamente ortogonales; sus normas al cuadrado
son $2L$, $L$ y $L$, respectivamente. Puede comprobarse con las
identidades de producto a suma, por ejemplo
$2\cos a\cos b=\cos(a-b)+\cos(a+b)$.
Una función impar integrable integra cero sobre $[-L,L]$.
Al comparar contribuciones a una aproximación importa el coeficiente
\emph{y} la norma de su función, no sólo el coeficiente.

%
\subsection*{Coeficientes, periodo y extensión}
\label{CH08-B02}
Para periodo $2L$, usamos la convención
\[
 f(x)\sim\frac{a_0}{2}+\sum_{n=1}^\infty
 \left[a_n\cos\frac{n\pi x}{L}+b_n\sin\frac{n\pi x}{L}\right],
\]
\[
 a_n=\frac1L\int_{-L}^{L}f(x)\cos\frac{n\pi x}{L}\,dx
 \ (n\ge0),\qquad
 b_n=\frac1L\int_{-L}^{L}f(x)\sin\frac{n\pi x}{L}\,dx.
\]
En una función par, $b_n=0$; en una impar, $a_n=0$. Para datos en
$(0,L)$, una extensión par da coeficientes de coseno mediante
$\frac2L\int_0^L f\cos(n\pi x/L)\,dx$; una impar da análogamente
coeficientes de seno. El observable físico decide la extensión: una
componente con signo y su módulo pueden tener simetrías distintas.
Compruébense periodo, frecuencia y unidades antes de integrar.
Consultas: soluciones de \exref{M2-E0051} y \exref{M2-E0345}.

Para interpretar un perfil que también depende del tiempo, una derivada
parcial mantiene fija la otra variable. En $u(x,t)=A(t)X(x)$,
$u_t=A'X$ y $u_{xx}=AX''$; sustituir en la ecuación dada permite obtener
una EDO de amplitud. Temperatura prescrita fija $u$; flujo nulo fija la
derivada espacial correspondiente cuando el flujo es proporcional a
$-u_x$. El desarrollo de estos modelos se recoge en el capítulo 10.

%
\subsection*{Suma de la serie y empalmes}
\label{CH08-B03}
Si la extensión periódica es suave a trozos, con límites laterales y un
número finito de piezas regulares por periodo, su serie de Fourier
converge en cada punto al promedio
\[
 \frac{f(x^-)+f(x^+)}{2}.
\]
En un punto de continuidad coincide con $f(x)$. En los extremos de un
periodo, los límites se leen usando el empalme periódico. Cambiar un
valor aislado no cambia los coeficientes. Una esquina continua no es un
salto del valor de la función.

La suma parcial hasta el armónico $N$ puede contener menos de $N$ términos
no nulos por simetría. Distíngase ese índice del número de componentes
retenidas. Cerca de un salto, el fenómeno de Gibbs produce oscilaciones
cuya zona se estrecha al aumentar $N$, sin desaparecer su sobreoscilación
relativa al salto. No se interprete una gráfica truncada como igualdad
puntual exacta.

%
\subsection*{Proyección y error cuadrático}
\label{CH08-B04}
Para una familia finita de funciones ortogonales no nulas $\phi_j$, la
mejor aproximación en el espacio que generan
respecto del error cuadrático usa
$c_j=\langle f,\phi_j\rangle/\|\phi_j\|^2$. En la convención anterior,
para $f$ de cuadrado integrable, Parseval da
\[
 \frac1L\int_{-L}^{L}|f|^2\,dx
 =\frac{a_0^2}{2}+\sum_{n\ge1}(a_n^2+b_n^2),\qquad
 \frac1L\int_{-L}^{L}|f-S_N|^2\,dx
 =\sum_{n>N}(a_n^2+b_n^2).
\]
Si $g$ es positiva y decreciente,
$\int_{N+1}^{\infty}g(x)\,dx\le\sum_{n>N}g(n)
\le\int_N^{\infty}g(x)\,dx$; compruébense esas hipótesis antes de
acotar una cola. La cota obtenida es integral, no uniforme ni puntual
cerca de un salto.

La media es $a_0/2$ y el promedio del cuadrado es
$a_0^2/4+\frac12\sum_{n\ge1}(a_n^2+b_n^2)$; su raíz es el valor
eficaz o media cuadrática. Un periodo elegido puede ser múltiplo del mínimo: revísense las
frecuencias realmente presentes. Consulta: solución de \exref{M2-E0348}.
